% GENERATED FILE. Edit canonical lessons, then run npm run generate:books.
% canonical-source-hash: fnv1a64:eae25b8aaa2eb27b
% canonical-lessons: KA-C147-kudi, KA-C147-odu, KA-C147-malagu, KA-C147-kattarisu, KA-C147-ode

\chapter{To drink, To run, To sleep, To cut, To break}
\label{ch:ka-to-drink-to-run-to-sleep-to-cut-to-break}
\hlchaptermodality{147}

\begin{chapteropening}
\textbf{By the end of this chapter:} \emph{I can say to drink, to run, to sleep, to cut and to break.}

Complete the last lesson of chapter 147: I can say to drink, to run, to sleep, to cut and to break.
\end{chapteropening}

\section[\texorpdfstring{\kn{ಕುಡಿ} (kuḍi)}{ಕುಡಿ (kuḍi)}]{\kn{ಕುಡಿ} (kuḍi) — to drink}

\label{lesson:KA-C147-kudi}



\begin{quote}
Before the new one: say the Kannada for golden, then the Kannada for silver.
\end{quote}

\subsection*{You'll want to know: \kn{ಕುಡಿ}}
\textbf{\kn{ಕುಡಿ}} — \emph{kuḍi} — \textquotedblleft{}to drink\textquotedblright{}.

\begin{cousinweb}
\textbf{\kn{ಕುಡಿ}} (\emph{kuḍi}) is native Dravidian, one word with Tamil \emph{kuṭi} and Malayalam \emph{kuṭi} (DEDR 1654). DEDR lists \emph{kuḍi} itself, \textquotedblleft{}to drink\textquotedblright{}. DEDR 1654 compares Sanskrit \emph{kuṭī}.
\end{cousinweb}

\begin{grammarlens}[title={What you've built}]
One more word for this chapter.
\end{grammarlens}

\subsection*{Your turn}
\begin{itemize}
\raggedright
  \item \emph{Say it:} \emph{kuḍi}
  \item \emph{Say it:} \emph{kuḍi}, once more
  \item \emph{Say it:} say \emph{kuḍi}
  \item \emph{Recall:} say the Kannada for golden, then the Kannada for silver, then say \emph{kuḍi} again
\end{itemize}

\begin{tcolorbox}[breakable,colback=teal!4,colframe=teal!35!black,title={Before you move on}]
What does \kn{ಕುಡಿ} mean? (\textquotedblleft{}To drink\textquotedblright{}.) Say it once more.
\end{tcolorbox}
\section[\texorpdfstring{\kn{ಓಡು} (ōḍu)}{ಓಡು (ōḍu)}]{\kn{ಓಡು} (ōḍu) — to run}

\label{lesson:KA-C147-odu}



\begin{quote}
Before the new one: say the Kannada for silver, then the Kannada for to drink.
\end{quote}

\subsection*{You'll want to know: \kn{ಓಡು}}
\textbf{\kn{ಓಡು}} — \emph{ōḍu} — \textquotedblleft{}to run\textquotedblright{}.

\begin{cousinweb}
\textbf{\kn{ಓಡು}} (\emph{ōḍu}) is native Dravidian, one word with Tamil \emph{ōṭu} and Telugu \emph{ōḍu} (DEDR 1041). DEDR lists \emph{ōḍu} itself, \textquotedblleft{}to run, flee\textquotedblright{}.
\end{cousinweb}

\begin{grammarlens}[title={What you've built}]
One more word for this chapter.
\end{grammarlens}

\subsection*{Your turn}
\begin{itemize}
\raggedright
  \item \emph{Say it:} \emph{ōḍu}
  \item \emph{Say it:} \emph{ōḍu}, once more
  \item \emph{Say it:} say \emph{ōḍu}
  \item \emph{Recall:} say the Kannada for silver, then the Kannada for to drink, then say \emph{ōḍu} again
\end{itemize}

\begin{tcolorbox}[breakable,colback=teal!4,colframe=teal!35!black,title={Before you move on}]
What does \kn{ಓಡು} mean? (\textquotedblleft{}To run\textquotedblright{}.) Say it once more.
\end{tcolorbox}
\section[\texorpdfstring{\kn{ಮಲಗು} (malagu)}{ಮಲಗು (malagu)}]{\kn{ಮಲಗು} (malagu) — to sleep, to lie down}

\label{lesson:KA-C147-malagu}



\begin{quote}
Before the new one: say the Kannada for to drink, then the Kannada for to run.
\end{quote}

\subsection*{You'll want to know: \kn{ಮಲಗು}}
\textbf{\kn{ಮಲಗು}} — \emph{malagu} — \textquotedblleft{}to sleep, to lie down\textquotedblright{}.

\begin{cousinweb}
\textbf{\kn{ಮಲಗು}} (\emph{malagu}) is native Dravidian, one word with Tulu \emph{malaguni}, \textquotedblleft{}to lie down\textquotedblright{} (DEDR 4735). DEDR lists \emph{malagu} itself, \textquotedblleft{}to recline, lie down\textquotedblright{}, and points to Tamil \emph{malar}, \textquotedblleft{}to lie on the back\textquotedblright{} (DEDR 4740).
\end{cousinweb}

\begin{grammarlens}[title={What you've built}]
One more word for this chapter.
\end{grammarlens}

\subsection*{Your turn}
\begin{itemize}
\raggedright
  \item \emph{Say it:} \emph{malagu}
  \item \emph{Say it:} \emph{malagu}, once more
  \item \emph{Say it:} say \emph{malagu}
  \item \emph{Recall:} say the Kannada for to drink, then the Kannada for to run, then say \emph{malagu} again
\end{itemize}

\begin{tcolorbox}[breakable,colback=teal!4,colframe=teal!35!black,title={Before you move on}]
What does \kn{ಮಲಗು} mean? (\textquotedblleft{}To sleep, to lie down\textquotedblright{}.) Say it once more.
\end{tcolorbox}
\section[\texorpdfstring{\kn{ಕತ್ತರಿಸು} (kattarisu)}{ಕತ್ತರಿಸು (kattarisu)}]{\kn{ಕತ್ತರಿಸು} (kattarisu) — to cut}

\label{lesson:KA-C147-kattarisu}



\begin{quote}
Before the new one: say the Kannada for to run, then the Kannada for to sleep.
\end{quote}

\subsection*{You'll want to know: \kn{ಕತ್ತರಿಸು}}
\textbf{\kn{ಕತ್ತರಿಸು}} — \emph{kattarisu} — \textquotedblleft{}to cut\textquotedblright{}.

\begin{cousinweb}
\textbf{\kn{ಕತ್ತರಿಸು}} (\emph{kattarisu}) is \emph{kattari}, \textquotedblleft{}scissors\textquotedblright{}, from Sanskrit \emph{kartarī} (CDIAL 2858), made into a verb with \emph{-isu}: \textquotedblleft{}to scissor\textquotedblright{}, to cut.
\end{cousinweb}

\begin{grammarlens}[title={What you've built}]
One more word for this chapter.
\end{grammarlens}

\subsection*{Your turn}
\begin{itemize}
\raggedright
  \item \emph{Say it:} \emph{kattarisu}
  \item \emph{Say it:} \emph{kattarisu}, once more
  \item \emph{Say it:} say \emph{kattarisu}
  \item \emph{Recall:} say the Kannada for to run, then the Kannada for to sleep, then say \emph{kattarisu} again
\end{itemize}

\begin{tcolorbox}[breakable,colback=teal!4,colframe=teal!35!black,title={Before you move on}]
What does \kn{ಕತ್ತರಿಸು} mean? (\textquotedblleft{}To cut\textquotedblright{}.) Say it once more.
\end{tcolorbox}
\section[\texorpdfstring{\kn{ಒಡೆ} (oḍe)}{ಒಡೆ (oḍe)}]{\kn{ಒಡೆ} (oḍe) — to break}

\label{lesson:KA-C147-ode}



\begin{quote}
Before the new one: say the Kannada for to sleep, then the Kannada for to cut.
\end{quote}

\subsection*{You'll want to know: \kn{ಒಡೆ}}
\textbf{\kn{ಒಡೆ}} — \emph{oḍe} — \textquotedblleft{}to break\textquotedblright{}.

\begin{cousinweb}
\textbf{\kn{ಒಡೆ}} (\emph{oḍe}) is native Dravidian, one word with Tamil \emph{oṭi} and Telugu \emph{oḍapi} (DEDR 946). DEDR lists Kannada \emph{oḍi}, \textquotedblleft{}to be broken, burst\textquotedblright{}.
\end{cousinweb}

\begin{grammarlens}[title={What you've built}]
One more word for this chapter.
\end{grammarlens}

\subsection*{Your turn}
\begin{itemize}
\raggedright
  \item \emph{Say it:} \emph{oḍe}
  \item \emph{Say it:} \emph{oḍe}, once more
  \item \emph{Say it:} say \emph{oḍe}
  \item \emph{Recall:} say the Kannada for to sleep, then the Kannada for to cut, then say \emph{oḍe} again
\end{itemize}

\begin{tcolorbox}[breakable,colback=teal!4,colframe=teal!35!black,title={Before you move on}]
What does \kn{ಒಡೆ} mean? (\textquotedblleft{}To break\textquotedblright{}.) Say it once more.
\end{tcolorbox}
